\section{The Soroban Boundary}\label{sec:soroban}

On chain, an agent's grant lives in a \emph{grant signer}
(\file{contracts/grant-signer}): a contract deployed once per account, holding
for each grant a rule-eval rule set (a root grant's rules, the ones the
account's scoped document for its agents lowers to, and for a delegated
grant the attenuation steps it adds), its holder's ed25519 key, its parent,
its budget node and a revoked flag. The account is the Perch account Nido's
factory mints, unchanged.

\subsection{A delegated signer, not a policy}

A policy contract attached to an agent's context rule is the usual way to add a
check to an OpenZeppelin smart account~\cite{ozStellarAccounts}. It would be
the wrong place here, for two reasons. Perch accounts are document-only, and
their document schema has no member that names another policy contract. More
importantly, a policy is asked only about the contexts the submitter routed to
its rule; every other context's rule would have to be shown unsatisfiable by an
agent. The document schema does declare delegated signers, which the account
authenticates through CAP-0071's \code{delegate\_account\_auth}, and the host
then calls the delegate's \code{\_\_check\_auth} with the account's payload
and its whole context list\claim{cap71-contexts}. So the account's document
names the grant signer as a signer and makes it the only signer of the rules
agents act through:
\begin{lstlisting}[basicstyle=\ttfamily\scriptsize]
"signers": [
  {"id": "owner", "verifier": "<WebAuthn verifier>",
   "key": "<passkey>"},
  {"id": "grants", "address": "<grant signer>"}],
"rules": [
  {"name": "admin", "scope": {"type": "self-admin"},
   "principals": {"type": "all", "signers": ["owner"]}},
  {"name": "a-token",
   "scope": {"type": "contract", "address": "<token>"},
   "principals": {"type": "all", "signers": ["grants"]}}]
\end{lstlisting}
The grant signer sees every context of the authorization, including contexts
the submitter routed to other rules, so the guarantee does not depend on rule
selection, which a delegate's signature does not bind. The payload a holder
signs is the host's payload for the account's entry, which names the account,
so a holder's signature authorizes for that one account. Adding agents to an
account is a document change and one deployment; nothing in Perch or
OpenZeppelin changes.

\subsection{Checking an authorization}

Figure~\ref{fig:soroban} follows one agent authorization. The account's
\code{\_\_check\_auth} runs Perch's guards and then OpenZeppelin's rule
check, which authenticates the rule's signer, the grant signer, by
\code{delegate\_account\_auth}. The grant signer's \code{\_\_check\_auth}
receives the payload, a signature naming the grant and carrying the holder's
ed25519 signature, and every context, and refuses at the first failure:
\begin{enumerate}
\item load the signing grant and its ancestors, refusing if any is revoked;
\item verify the holder's signature over the payload;
\item refuse a contract creation; a context that calls the grant signer
  itself, unless it is the holder's own delegation; a context that calls the
  account; or a call on the budget asset other than \code{transfer} with a
  non-negative amount;
\item read every other context through rule-core's Soroban facade and
  require rule-eval's \code{authorize} over the grant's rule set, each
  context by the first rule that allows it, with the grant's cap records;
\item charge each budget-asset transfer to the grant and every ancestor with
  rule-eval's \code{spend}, refusing if any node would be overdrawn, and only then
  store the new \code{spent} values.
\end{enumerate}
The grant signer writes nothing before it refuses. A refusal fails the
account's authorization, and with it the call that asked for it; a calling
contract can catch that failure and go on, so the transaction need not fail.
Soroban then rolls back the failed call's storage writes and authorization
state, so either way the refusal changes nothing. The holder's ed25519 signature covers the host's payload,
not the grant id carried beside it, so a key that holds two grants produces
signatures usable under either; one key per grant avoids this.

\paragraph{What the grant signer sees.} The contexts it checks are those that
need the account's authorization. A call's effects that need none are outside
the tree and outside every caveat and budget: a \code{transfer\_from} against
an allowance the owner gave a contract earlier, a contract that already holds
or controls the account's funds, or any allowed call whose side effects never
call \code{require\_auth} for the account. The guarantee against hidden calls
holds for calls that need the account's authorization, and the assumption that
no allowed contract has such a standing position (\S\ref{sec:threat}) carries
the rest\claim{no-standing-allowance}.

\begin{figure}[t]
\centering
\resizebox{\linewidth}{!}{\input{soroban-flow}}
\caption{An agent's call through the account. The host passes every context of
the account's authorization to the grant signer, which checks them all with
the evaluator; one disallowed context refuses the whole tree.}
\label{fig:soroban}
\Description{An authorization tree in which an allowed router call nests a
token transfer; both contexts reach the grant signer, which runs rule-eval's authorize.}
\end{figure}

\paragraph{Delegation, and the grant signer's second role.}
\code{delegate(parent, key, steps, allowance)} stores a child grant whose
rule set is the parent's with the given rule-eval attenuation steps applied,
with a budget record under the parent's. It requires the grant signer's own authorization, so the
host calls the grant signer's \code{\_\_check\_auth} in two roles, with the
same three arguments: as the account's delegate, with the account's payload
and contexts, and as the check of the grant signer's own address, whenever a
call needs that address's authorization (its own \code{delegate}, a call on
another contract that asks for it, a deployment with it as deployer). The
host lets a contract's \code{require\_auth} of itself call its own
\code{\_\_check\_auth} and refuses any other re-entry. The checks are the
same in both roles: contract creation is refused under every
grant\claim{signer-deploy}, since in the second role no account document
stands behind the check; a call on the account, an unmetered call on the
budget asset and a call on the grant signer other than the holder's own
delegation are refused; and the rule set bounds every other call. A
signature in the second role authorizes nothing for the account, which
accepts the grant signer only through
\code{delegate\_account\_auth}\claim{signer-own-address}. One consequence is
stated rather than enforced. Any contract can call \code{require\_auth} on an
address it is given, so a holder can have a contract its grant allows ask
for the grant signer's own authorization, and authorize what that contract
does as the grant signer's address, within the grant's rules. That includes
moving what the address holds; a transfer of the budget asset is still
charged to the grant, and one of any other asset is bounded only by the
rules.
The deployment gives the address no balance: a contract's balance in a
Stellar Asset Contract, the native one included, is absent until something
is sent to it, and storage rent is paid in fees by the transaction's
submitter. But the address is published in the account's document, and
whatever is sent to it is reachable this way.

\subsection{Threat coverage}

Table~\ref{tab:soroban-threats} maps each attempt in \S\ref{sec:threat} to
what refuses it. Each refusal is asserted in an integration test by the
contract that refused and its error, on a factory-minted Perch account under
enforcing authorization, where the host runs every \code{\_\_check\_auth}.

\begin{table*}[t]
\caption{Escalation attempts at the Soroban boundary and what refuses them.}
\label{tab:soroban-threats}
\small
\begin{tabularx}{\linewidth}{@{}>{\raggedright\arraybackslash}p{0.27\linewidth}>{\raggedright\arraybackslash}X l@{}}
\toprule
\textbf{Attempt} & \textbf{What refuses it} & \textbf{Status} \\
\midrule
Issue, revoke, rotate or refill a grant, its own included &
  These need the account's authorization, which only the owner's admin rule
  gives; no agent rule is scoped to the grant signer, and asked directly the
  grant signer refuses every call on itself but the holder's own delegation. &
  \claim{signer-admin} \\
Change the account's rules or policies, or call its admin functions &
  No agent rule is scoped to the account; a document that wrongly scopes one
  still meets the grant signer's refusal of calls on the account. &
  \claim{signer-account} \\
Delegate from an ancestor's grant, or one it does not hold &
  The grant signer accepts a delegation only with the signing grant as parent. &
  \claim{signer-delegate-other} \\
Widen a grant by delegating &
  A child's rule set is its parent's with rule-eval attenuation steps
  applied, each of which narrows. &
  \claim{signer-narrows}\ \claim{signer-rule-eval-chain} \\
Hide a disallowed call inside an allowed one &
  Every context of the account's authorization is checked; a nested call the
  signed tree leaves out fails its \code{require\_auth} in the host. Effects
  that need no authorization from the account are outside the tree. &
  \claim{signer-hidden-call}\ \claim{batch-narrower} \\
Deploy a contract &
  The grant signer refuses every contract creation, under every grant and in
  both roles; through the account, a Perch document has no creation scope,
  so the account refuses it first. &
  \claim{signer-deploy}\ \claim{create-contract-denied} \\
Hand off the budget asset (\code{approve}, \code{burn}) &
  Only metered transfers of the budget asset pass. &
  \claim{signer-metered} \\
Overspend through siblings or several transfers &
  Every transfer in the authorization is charged to the spender and every
  ancestor; the first overdrawn node refuses and nothing is stored. &
  \claim{signer-oversubscribed}\ \claim{signer-all-charged}\ \claim{signer-budget-total} \\
Keep acting after revocation, or with a rotated-out key &
  Revocation of any ancestor refuses; rotation replaces the key. &
  \claim{signer-revocation}\ \claim{signer-rotation} \\
\bottomrule
\end{tabularx}
\end{table*}

\subsection{The wiring}

The evaluator is verified; what connects it to Soroban is not. The grant
signer reads each context through rule-core's Soroban facade, as the relay
reads a call's XDR: a contract and an address argument by its
\code{ScAddress} XDR, a function and a symbol by its bytes, \code{i128},
\code{u64}, \code{bool}, bytes and strings by value, and any other argument
as \code{other} with its full XDR. A test checks that the grant signer and
the relay read every argument type alike\claim{signer-relay-mapping}; that
the reading keeps distinct arguments distinct rests on XDR encoding being
injective and is not proven\claim{xdr-injective}. On every grant that names
its contracts, the rule set a caveat list lowers to allows exactly what the
caveats allow\claim{chain-rule-eval-agree}, and the grant signer reproduces a
recorded corpus of the outcomes it had when it held caveat lists and checked
them with the chain evaluator\claim{signer-check-auth-corpus}.

\paragraph{The budget link.} The grant signer never holds the ledger of
every grant whole. Its storage holds one grant per id, whose parent,
allowance and \code{spent} are that grant's node; \code{issue} and
\code{delegate} store a grant at the next id with nothing spent, which is the
model's \code{append}; \code{refill} calls rule-eval's \code{refill} on a
ledger of the one grant; and each \code{\_\_check\_auth} builds a ledger of only
the signing grant's lineage, root first with each node's parent the one before
it, calls rule-eval's \code{spend} on it once per budget-asset transfer, stopping
at the first failure, and only if every spend succeeded writes the changed
\code{spent} values back. \file{formal/AttenuatingGrants/Signer.lean} follows
these steps. Spending on the lineage ledger has the whole ledger's outcome at
the signing grant, an error naming the same grant, and afterwards the whole
ledger's lineage, and the spend loop over grant-eval's extracted \code{spend}
computes it\claim{lineage-ledger}. The contract's rule-eval \code{spend}
agrees with that extracted \code{spend} on lineage ledgers, on every outcome,
overdrawn grant and ledger, for amounts up to the largest \code{i128}, which
\code{issue}, \code{delegate} and \code{refill} keep every allowance
within\claim{signer-rule-eval-budget}. One authorization's spends are all or nothing:
the contract stores them only after every spend succeeds, and it does so in
the \code{require\_auth} of the authorization's root, inside the call that
encloses every call the authorization allows. When that call fails, whether
or not a caller catches the failure, the host rolls back its storage writes,
the charge included, and a charge stays only with the root call it was made
in. A test catches a refused authorization and an authorization whose
nested transfer fails after the charge, and finds the budget unchanged both
times\claim{host-frame-rollback}. So every sequence of the grant signer's calls is a run of the model on the
ledger of every grant\claim{budget-batch-as-run}, and Theorem~\ref{thm:budget}'s
bound holds for it: the total spent at or under any grant never exceeds that
grant's allowance, with no hypothesis on the
allowances\claim{signer-budget-total}. What remains trusted is that the
contract's few lines of storage, lineage and loops are those
definitions\claim{wiring-budget-calls}. The tests check the outcome:
oversubscribed siblings cannot together pass their parent, and a refused
spend stores nothing\claim{signer-oversubscribed}.

Two requirements on the account's document are stated and met by the
documents the tests apply, but nothing on chain checks them: a holder key must
reach the account only through the grant signer, since registering it as a
signer by another route lets it authorize without its grant; and the document
must name its own account's grant signer, since the grant signer cannot see
which account delegated to it\claim{document-requirements}.
